\documentclass[10pt,a4paper]{amsart}
\usepackage[margin=19mm]{geometry}
\usepackage{amsmath,amssymb,mathtools}
\usepackage[scr=rsfs,cal=euler]{mathalfa}
\usepackage{xfrac}
\usepackage[unicode,hidelinks,bookmarks=false]{hyperref}
\newcommand{\ie}{i.e.\ }
\newcommand{\cf}{cf.\ }
\newcommand{\resp}{respectively\ }
\newcommand{\Subsection}{\subsection}
\providecommand{\nobreakdash}{\nobreak\hbox{-}}
\setlength{\emergencystretch}{2em}
\multlinegap0pt
\usepackage{bm}
\usepackage{bbm}
\usepackage{dsfont}
\usepackage{xspace}
\usepackage{cancel}
\usepackage{mathtools}
\usepackage{amsthm}
\makeatletter
\@ifundefined{lemma}{\newtheorem{lemma}{Lemma}}{}
\@ifundefined{proposition}{\newtheorem{proposition}{Proposition}}{}
\@ifundefined{theorem}{\newtheorem{theorem}{Theorem}}{}
\makeatother
\title[Robust and Smooth Estimation of the Extreme Tail Index via W]{Robust and Smooth Estimation of the Extreme Tail Index via Weighted Minimum Density Power Divergence}
\author{Saida Mancer, Abdelhakim Necir, Djamel Meraghni}
\date{Source: arXiv:2507.15744v2 (CC BY 4.0)}
\begin{document}
\maketitle

\noindent\emph{Source and licence.} Robust and Smooth Estimation of the Extreme Tail Index via Weighted Minimum Density Power Divergence. Version arXiv:2507.15744v2 (CC BY 4.0), distributed under the Creative Commons Attribution 4.0 International licence, as recorded in the arXiv licence metadata for this exact version and shown on the abstract page https://arxiv.org/abs/2507.15744v2. Full rights notice: licenses/arxiv-2507.15744-CC-BY-4.0.txt.\par\smallskip

\noindent\textit{Editorial scope.} Counted unit: the Lemma whose source label is \texttt{lemma4} in the article (source0.tex lines 1862--1867, source label \texttt{lemma4}), delivered together with the article's own complete original proof of it (source0.tex lines 1869--1984, one \texttt{proof} environment of 116 lines). No mathematics is written, repaired or re-derived: statement and proof are the article's own text line for line. Required context retained uncounted: RV (lines 98--101); L (lines 495--505); second-order (lines 594--598); Theorem2 (lines 607--660); psi-m (lines 757--761); lemma2 (lines 1378--1392); eta (lines 1536--1540); prop1 (lines 1986--2009); prop2 (lines 2110--2133). Every definition, display and label of those lines is retained unchanged. Omitted from this package and not redistributed: 1--97, 102--494, 506--593, 599--606, 661--756, 762--1377, 1393--1535, 1541--1861, 1868--1868, 1985--1985, 2010--2109, 2134--2716. Nothing inside a retained excerpt is omitted or altered: every retained body is delivered exactly as the article's lines are written, so no omission receipt of any kind accompanies this package. Citations: the retained excerpts carry no citation command at all, so no bibliography entry of the article is transcribed and no reference is redistributed. Reference adaptation: none was needed and none was made; every reference inside the delivered unit resolves inside it, so no unresolved reference or citation is produced. Printed slips kept literal: no printed word, symbol, hypothesis or number of the counted statement or of its proof was altered. Layout and class: the article's own class is \texttt{amsart} with packages including lineno, natbib, hyperref, caption, amsfonts, amsmath, amssymb, graphicx; the wrapper is the lane's pinned \texttt{amsart} editorial wrapper at 10pt on A4, which restates the theorem-like environments the retained text needs and adds the article's own macro definitions that the retained text needs (none), with extra wrapper packages (bm, bbm, dsfont, xspace, cancel, mathtools). The delivered numbering is the unit's own. Assets: no retained span contains a figure environment, a graphics-inclusion command, a PDF-page inclusion, a table or a TikZ picture; no image file is copied into this package, the package declares no asset dependency and no image file is redistributed with it. Licence: version arXiv:2507.15744v2 of this article is distributed under the Creative Commons Attribution 4.0 International licence, as recorded in the arXiv licence metadata for this exact version and shown on the abstract page https://arxiv.org/abs/2507.15744v2. A full rights notice is delivered at licenses/arxiv-2507.15744-CC-BY-4.0.txt. Characters: every retained excerpt is pure ASCII, so the delivered engine, XeLaTeX with its default Latin Modern text font, reports no missing character.
\begin{equation}
\frac{\overline{F}\left(  ux\right)  }{\overline{F}\left(  u\right)
}\rightarrow x^{-1/\gamma},\text{ as }u\rightarrow\infty. \label{RV}%
\end{equation}

\begin{equation}
\dfrac{1}{k}%
%TCIMACRO{\dsum \limits_{i=1}^{k}}%
%BeginExpansion
{\displaystyle\sum\limits_{i=1}^{k}}
%EndExpansion
J\left(  \dfrac{i}{k+1}\right)  \dfrac{d}{d\gamma}\ell_{\gamma,J}^{\alpha
}\left(  \dfrac{X_{n-i+1:n}}{X_{n-k:n}}\right)  =\dfrac{\alpha}{\alpha+1}%
\int_{1}^{\infty}\dfrac{d}{d\gamma}\ell_{\gamma,J}^{1+\alpha}\left(  x\right)
dx,\text{ for }\alpha>0. \label{L}%
\end{equation}

\begin{equation}
\lim_{t\rightarrow\infty}\frac{\overline{F}\left(  tx\right)  /\overline
{F}\left(  t\right)  -x^{-1/\gamma_{0}}}{A\left(  t\right)  }=x^{-1/\gamma
_{0}}\dfrac{x^{\tau/\gamma_{0}}-1}{\tau\gamma_{0}}, \label{second-order}%
\end{equation}

\begin{theorem}
\textbf{\label{Theorem2}}(asymptotic normality) Assume that tail distribution
$\overline{F}$ satisfies the second-order condition $\left(
\ref{second-order}\right)  $ and that, for $\alpha>0,$ the estimator
$\widehat{\gamma}_{k,\alpha,J}$ is a consistent for $\gamma_{0}$ satisfying
$\left(  \ref{L}\right)  .$ Let $J$ be a continuous function fulfilling
assumptions $\left[  A1\right]  -\left[  A3\right]  $ and $k$ be a sequence of
integers such that $k\rightarrow\infty,$ $k/n\rightarrow0$ and $\sqrt
{k}a\left(  n/k\right)  $ is asymptotically bounded. Then, in the probability
space $\left(  \Omega,\mathcal{A},\mathbf{P}\right)  $ there exists a standard
Wiener process $\left\{  W\left(  x\right)  ,\text{ }x\geq0\right\}  ,$ such
that
\begin{align*}
&  \left(  1+\frac{1}{\alpha}\right)  \eta_{\gamma_{0},\alpha}\sqrt{k}\left(
\widehat{\gamma}_{k,\alpha,J}-\gamma_{0}\right) \\
&  =\int_{1}^{\infty}\left(  W\left(  x^{-1/\gamma_{0}}\right)  -x^{-1/\gamma
_{0}}W\left(  1\right)  \right)  J\left(  x^{-1/\gamma_{0}}\right)
d\Psi_{\gamma_{0},\alpha}^{\left(  1\right)  }+\sqrt{k}a\left(  n/k\right)
B_{\gamma_{0},\alpha}^{\left(  1\right)  }+o_{\mathbf{P}}\left(  1\right)  ,
\end{align*}
where $\Psi_{\gamma_{0},\alpha}^{\left(  1\right)  }$ and $\eta_{\gamma
_{0},\alpha}$ are respectively given in $\left(  \ref{psi-m}\right)  $ and
$\left(  \ref{eta}\right)  ,$ and
\[
B_{\gamma_{0},\alpha}^{\left(  1\right)  }:=\int_{1}^{\infty}x^{-1/\gamma_{0}%
}\frac{x^{\tau/\gamma_{0}}-1}{\tau\gamma_{0}}J\left(  x^{-1/\gamma_{0}%
}\right)  d\Psi_{\gamma_{0},\alpha}^{\left(  1\right)  }\left(  x^{-1/\gamma
_{0}}\right)  ,
\]
Whenever $\sqrt{k}a\left(  n/k\right)  \rightarrow\lambda<\infty,$
\[
\left(  1+\frac{1}{\alpha}\right)  \eta_{\gamma_{0},\alpha}\sqrt{k}\left(
\widehat{\gamma}_{k,\alpha,J}-\gamma_{0}\right)  \overset{\mathcal{D}%
}{\rightarrow}\mathcal{N}\left(  \lambda B_{\gamma_{0},\alpha}^{\left(
1\right)  },\sigma_{\gamma_{0},\alpha}^{2}\right)  ,
\]
as $n\rightarrow\infty,$ where%
\[
\sigma_{\gamma_{0},\alpha}^{2}:=%
%TCIMACRO{\dint _{0}^{1}}%
%BeginExpansion
{\displaystyle\int_{0}^{1}}
%EndExpansion%
%TCIMACRO{\dint _{0}^{1}}%
%BeginExpansion
{\displaystyle\int_{0}^{1}}
%EndExpansion
\left(  \min\left(  s,t\right)  -st\right)  J\left(  s\right)  J\left(
t\right)  d\Psi_{\gamma_{0},\alpha}^{\left(  1\right)  }\left(  s^{-\gamma
_{0}}\right)  d\Psi_{\gamma_{0},\alpha}^{\left(  1\right)  }\left(
t^{-\gamma_{0}}\right)  .
\]

\end{theorem}

\begin{equation}
\Psi_{\gamma_{0},\alpha}^{\left(  m\right)  }\left(  x\right)  :=\left.
\frac{d^{m}\ell_{\gamma,J}^{\alpha}\left(  x\right)  }{d^{m}\gamma}\right\vert
_{\gamma=\gamma_{0}}, \label{psi-m}%
\end{equation}

\begin{lemma}
\textbf{\label{lemma2}}Assume that $\overline{F}$ satisfies the first-order
condition $\left(  \ref{RV}\right)  .$ Let $J$ be a continuous function
fulfilling assumptions $\left[  A1\right]  -\left[  A3\right]  $ and let $k$
be a sequence of integers such that $k\rightarrow\infty$ and $k/n\rightarrow
0,$ then
\[
\pi_{k}^{\left(  1\right)  }\left(  \gamma_{0}\right)  \overset{\mathbf{P}%
}{\rightarrow}0\text{ and \ }\pi_{k}^{\left(  2\right)  }\left(  \gamma
_{0}\right)  \overset{\mathbf{P}}{\rightarrow}\eta_{\gamma_{0},\alpha},\text{
as }n\rightarrow\infty,
\]
for every $\alpha>0,$ where $\eta_{\gamma_{0},\alpha}$ is as in Theorem
$\ref{Theorem2}.$
\end{lemma}

\begin{equation}
\eta_{\gamma_{0},\alpha}:=\left(  1+\alpha\right)  \int_{1}^{\infty}\left(
\Psi_{\gamma_{0},1}^{\left(  1\right)  }\left(  x\right)  \right)  ^{2}%
\ell_{\gamma_{0},J}^{\alpha-1}\left(  x\right)  dx, \label{eta}%
\end{equation}%

\begin{proposition}
\textbf{\label{prop1}}For every $\alpha>0,$ we have%
\[
\Psi_{\gamma,\alpha+1}^{\left(  1\right)  }\left(  x\right)  =\left(
1+\frac{1}{\alpha}\right)  \ell_{\gamma,J}\left(  x\right)  \Psi
_{\gamma,\alpha}^{\left(  1\right)  }\left(  x\right)  ,
\]%
\[
\Psi_{\gamma,\alpha+1}^{\left(  2\right)  }\left(  x\right)  =\left(
1+\frac{1}{\alpha}\right)  \ell_{\gamma,J}\left(  x\right)  \Psi
_{\gamma,\alpha}^{\left(  2\right)  }\left(  x\right)  +\left(  1+\alpha
\right)  \left(  \Psi_{\gamma,1}^{\left(  1\right)  }\left(  x\right)
\right)  ^{2}\ell_{\gamma,J}^{\alpha-1}\left(  x\right)
\]
and%
\[
\Psi_{\gamma,\alpha+1}^{\left(  3\right)  }\left(  x\right)  =\left(
1+\frac{1}{\alpha}\right)  \ell_{\gamma,J}\left(  x\right)  \Psi
_{\gamma,\alpha}^{\left(  3\right)  }\left(  x\right)  +g_{\gamma,\alpha
}\left(  x\right)  ,\text{ }%
\]
for some integrable function $x\rightarrow g_{\gamma,\alpha}\left(  x\right)
.$
\end{proposition}

\begin{proposition}
\textbf{\label{prop2}}Assume that assumption $\left[  A1\right]  -\left[
A3\right]  $ hold, then for every $\alpha>0,$ each derivative function
$\Psi_{\gamma,\alpha}^{\left(  m\right)  },$ $1\leq m\leq4,$ is a finite
linear combination of
\begin{equation}
E_{m,\alpha}^{\left(  i\right)  }\left(  x;\gamma\right)  :=\left(  x^{-a_{i}%
}\log^{b_{i}}x\right)  \ell_{\gamma,J}^{\alpha}\left(  x\right)
%TCIMACRO{\tprod _{j=1}^{m}}%
%BeginExpansion
{\textstyle\prod_{j=1}^{m}}
%EndExpansion
\left[  \mathcal{L}^{\left(  j-1\right)  }\left(  x^{-1/\gamma}\right)
\right]  ^{c_{i,j}}, \label{lc}%
\end{equation}
for some constant $a_{i}\geq0$ and integers $b_{i},c_{i,j}\geq0,$ where
$\mathcal{L}^{\left(  j\right)  }$ denotes the $j$-th derivative of
$\mathcal{L} $ and $\mathcal{L}^{\left(  0\right)  }:=\mathcal{L}.$ Moreover
$\Psi_{\gamma,\alpha}^{\left(  m\right)  }$ is bounded over $x>1,$ and both
$\int_{1}^{\infty}\left\vert \Psi_{\gamma,\alpha+1}^{\left(  m\right)
}\left(  x\right)  \right\vert dx$ and $\int_{1}^{\infty}\left\vert
\ell_{\gamma,J}\left(  x\right)  \Psi_{\gamma,\alpha}^{\left(  m\right)
}\left(  x\right)  \right\vert dx$ are finite, for $1\leq m\leq4.$
\end{proposition}

\begin{lemma}
\textbf{\label{lemma4}}Assume that assumptions $\left[  A1\right]  -\left[
A3\right]  ,$ then fo given a consistent estimator $\widehat{\gamma}$ of
$\gamma_{0},$ we have $\pi_{k}^{\left(  3\right)  }\left(  \widehat{\gamma
}\right)  =O_{\mathbf{P}}\left(  1\right)  ,$ provided that $\alpha>0.$
\end{lemma}

\begin{proof}
Let us first show that $\pi_{k}^{\left(  3\right)  }\left(  \gamma_{0}\right)
=O_{\mathbf{P}}\left(  1\right)  $ and recall that
\[
\pi_{k}^{\left(  3\right)  }\left(  \gamma_{0}\right)  =\int_{1}^{\infty}%
\Psi_{\gamma_{0},\alpha+1}^{\left(  3\right)  }\left(  x\right)  dx-\left(
1+\frac{1}{\alpha}\right)  \frac{1}{k}\sum_{i=1}^{k}J\left(  \frac{i}%
{k}\right)  \Psi_{\gamma_{0},\alpha}^{\left(  3\right)  }\left(
\frac{X_{n-i+1:n}}{X_{n-k:n}}\right)  .
\]
Making use of Proposition $\ref{prop1},$ we may rewrite $\pi_{k}^{\left(
3\right)  }\left(  \gamma_{0}\right)  $ into the sum of $\int_{1}^{\infty
}g_{\gamma_{0},\alpha}\left(  x\right)  dx$ and $\left(  1+\frac{1}{\alpha
}\right)  A_{k}^{\left(  3\right)  }\left(  \gamma_{0}\right)  ,$ where
\[
A_{k}^{\left(  3\right)  }\left(  \gamma_{0}\right)  :=\int_{1}^{\infty}%
\dfrac{d}{d\gamma}\ell_{\gamma_{0},J}\left(  x\right)  \Psi_{\gamma_{0}%
,\alpha}^{\left(  3\right)  }\left(  x\right)  dx-\frac{1}{k}\sum_{i=1}%
^{k}J\left(  \frac{i}{k}\right)  \Psi_{\gamma_{0},\alpha}^{\left(  3\right)
}\left(  \frac{X_{n-i+1:n}}{X_{n-k:n}}\right)  .
\]
By similar arguments as those used in the proof of Lemma $\ref{lemma2},$ we
also show that $A_{k}^{\left(  3\right)  }\left(  \gamma_{0}\right)
\overset{\mathbf{P}}{\rightarrow}0,$ so $\pi_{k}^{\left(  3\right)  }\left(
\gamma_{0}\right)  \overset{\mathbf{P}}{\rightarrow}\int_{1}^{\infty}%
g_{\gamma_{0},\alpha}\left(  x\right)  dx<\infty.$ From the third assertion of
Proposition $\ref{prop1},$ we have%
\[
g_{\gamma_{0},\alpha}\left(  x\right)  =\Psi_{\gamma_{0},\alpha+1}^{\left(
3\right)  }\left(  x\right)  -\left(  1+\frac{1}{\alpha}\right)  \ell
_{\gamma,J}\left(  x\right)  \Psi_{\gamma_{0},\alpha}^{\left(  3\right)
}\left(  x\right)  .
\]
By applying Proposition $\ref{prop2},$ we infer that $\int_{1}^{\infty
}\left\vert g_{\gamma_{0},\alpha}\left(  x\right)  \right\vert dx<\infty,$
thus $\pi_{k}^{\left(  3\right)  }\left(  \gamma_{0}\right)  =O_{\mathbf{P}%
}\left(  1\right)  .$ Next we prove that $\pi_{k}^{\left(  3\right)  }\left(
\widehat{\gamma}_{0}\right)  -\pi_{k}^{\left(  3\right)  }\left(  \gamma
_{0}\right)  =o_{\mathbf{P}}\left(  1\right)  ,$ as $n\rightarrow\infty.$
Indeed, let us write
\begin{align*}
&  \pi_{k}^{\left(  3\right)  }\left(  \widehat{\gamma}\right)  -\pi
_{k}^{\left(  3\right)  }\left(  \gamma_{0}\right)  \medskip\\
&  =\int_{1}^{\infty}\left\{  \Psi_{\widehat{\gamma},\alpha+1}^{\left(
3\right)  }\left(  x\right)  -\Psi_{\gamma_{0},\alpha+1}^{\left(  3\right)
}\left(  x\right)  \right\}  dx\medskip\\
&  -\left(  1+\frac{1}{\alpha}\right)  \frac{1}{k}\sum_{i=1}^{k}J\left(
\frac{i}{k}\right)  \left\{  \Psi_{\widehat{\gamma},\alpha}^{\left(  3\right)
}\left(  \frac{X_{n-i+1:n}}{X_{n-k:n}}\right)  -\Psi_{\gamma_{0},\alpha
}^{\left(  3\right)  }\left(  \frac{X_{n-i+1:n}}{X_{n-k:n}}\right)  \right\}
.
\end{align*}
Using the mean value theorem, yields%
\[
\int_{1}^{\infty}\left\{  \Psi_{\widehat{\gamma},\alpha+1}^{\left(  3\right)
}\left(  x\right)  -\Psi_{\gamma_{0},\alpha+1}^{\left(  3\right)  }\left(
x\right)  \right\}  dx=\left(  \widehat{\gamma}-\gamma_{0}\right)
\Psi_{\widehat{\gamma}_{0}^{\ast},\alpha+1}^{\left(  4\right)  }\left(
x\right)  \int_{1}^{\infty}\Psi_{\widehat{\gamma}_{0}^{\ast},\alpha
+1}^{\left(  4\right)  }\left(  x\right)  dx,
\]
where $\widehat{\gamma}_{0}^{\ast}$ is between $\widehat{\gamma}$ and
$\gamma.$ Recall that $\widehat{\gamma}\overset{\mathbf{P}}{\rightarrow}%
\gamma_{0}$ and by Proposition $\ref{prop2}$
\[
\int_{1}^{\infty}\left\vert \Psi_{\widehat{\gamma}_{0}^{\ast},\alpha
+1}^{\left(  4\right)  }\left(  x\right)  \right\vert dx=O_{\mathbf{P}}\left(
1\right)  ,
\]
thus $\int_{1}^{\infty}\left\{  \Psi_{\widehat{\gamma},\alpha+1}^{\left(
3\right)  }\left(  x\right)  -\Psi_{\gamma_{0},\alpha+1}^{\left(  3\right)
}\left(  x\right)  \right\}  dx=o_{\mathbf{P}}\left(  1\right)  .$ On the
other hand
\[%
\begin{array}
[c]{cl}
& \dfrac{1}{k}%
%TCIMACRO{\dsum \limits_{i=1}^{k}}%
%BeginExpansion
{\displaystyle\sum\limits_{i=1}^{k}}
%EndExpansion
J\left(  \dfrac{i}{k}\right)  \left\{  \Psi_{\widehat{\gamma},\alpha}^{\left(
3\right)  }\left(  \dfrac{X_{n-i+1:n}}{X_{n-k:n}}\right)  -\Psi_{\gamma
_{0},\alpha}^{\left(  3\right)  }\left(  \dfrac{X_{n-i+1:n}}{X_{n-k:n}%
}\right)  \right\}  \bigskip\\
= &
%TCIMACRO{\dint _{1}^{\infty}}%
%BeginExpansion
{\displaystyle\int_{1}^{\infty}}
%EndExpansion
J\left(  \dfrac{\overline{F}_{n}\left(  xX_{n-k:n}\right)  }{\overline{F}%
_{n}\left(  X_{n-k:n}\right)  }\right)  \left\{  \Psi_{\widehat{\gamma}%
,\alpha}^{\left(  3\right)  }\left(  x\right)  -\Psi_{\gamma_{0},\alpha
}^{\left(  3\right)  }\left(  x\right)  \right\}  d\dfrac{F_{n}\left(
xX_{n-k:n}\right)  }{\overline{F}_{n}\left(  X_{n-k:n}\right)  }.
\end{array}
\]
Once again making use of the mean value theorem, we write%
\[
\left(  \widehat{\gamma}-\gamma_{0}\right)  \int_{1}^{\infty}J\left(
\frac{\overline{F}_{n}\left(  xX_{n-k:n}\right)  }{\overline{F}_{n}\left(
X_{n-k:n}\right)  }\right)  \Psi_{\overline{\gamma}_{0}^{\ast},\alpha
}^{\left(  4\right)  }\left(  x\right)  d\frac{F_{n}\left(  xX_{n-k:n}\right)
}{\overline{F}_{n}\left(  X_{n-k:n}\right)  },
\]
where $\overline{\gamma}_{0}^{\ast}$ is between $\widehat{\gamma}$ and
$\gamma_{0}.$ From Proposition $\ref{prop2},$ the function $x\rightarrow
\Psi_{\overline{\gamma}_{0}^{\ast},\alpha}^{\left(  4\right)  }\left(
x\right)  $ being bounded over $x>1,$ then the previous quantity equals
$o_{\mathbf{P}}\left(  1\right)  \int_{0}^{1}J\left(  s\right)
ds=o_{\mathbf{P}}\left(  1\right)  .$ Thus we showed that $\pi_{k}^{\left(
3\right)  }\left(  \widehat{\gamma}\right)  -\pi_{k}^{\left(  3\right)
}\left(  \gamma_{0}\right)  =o_{\mathbf{P}}\left(  1\right)  ,$ leading to
$\pi_{k}^{\left(  3\right)  }\left(  \widehat{\gamma}\right)  =O_{\mathbf{P}%
}\left(  1\right)  $ as well.
\end{proof}

\end{document}
